% ================================================================
%  LEGAL LUMINAIRE — ठेकेदार हेमराज का बचाव संक्षेप (हिन्दी)
%  राजस्थान | निर्माण विधि | FSL रिपोर्ट को चुनौती
%  पाँच विशेषज्ञ एजेंटों द्वारा सत्यापित
%  संकलन विधि: xelatex defence_hindi.tex
%  अस्वीकरण: न्यायालय में उपयोग से पूर्व सभी IS उद्धरणों को BIS की
%  मूल प्रकाशनों से सत्यापित करें।
% ================================================================

\documentclass[12pt, a4paper]{article}

% --- XeLaTeX / Unicode packages ---
\usepackage{fontspec}
\usepackage{polyglossia}
\setmainlanguage{hindi}
\setotherlanguage{english}
\setmainfont{Noto Serif Devanagari}[Script=Devanagari, Scale=1.0]
\newfontfamily\devanagarifont{Noto Serif Devanagari}[Script=Devanagari]
\newfontfamily\englishfont{Liberation Serif}

% --- Layout packages ---
\usepackage[a4paper, margin=2.5cm, top=3cm, bottom=3cm]{geometry}
\usepackage{microtype}
\usepackage{setspace}
\usepackage{booktabs}
\usepackage{longtable}
\usepackage{array}
\\usepackage{tabularx}
\usepackage{xcolor}
\usepackage{mdframed}
\usepackage{titlesec}
\usepackage{fancyhdr}
\usepackage{hyperref}
\usepackage{enumitem}
\usepackage{parskip}

% --- Colours ---
\definecolor{navy}{RGB}{24, 42, 80}
\definecolor{gold}{RGB}{180, 140, 30}
\definecolor{lightgold}{RGB}{255, 248, 220}
\definecolor{lightred}{RGB}{255, 240, 240}
\definecolor{lightblue}{RGB}{240, 245, 255}
\definecolor{lightgreen}{RGB}{240, 255, 245}
\definecolor{midgrey}{RGB}{100, 100, 100}
\definecolor{darkgrey}{RGB}{50, 50, 50}

% --- Section formatting ---
\titleformat{\section}[block]{\large\bfseries\color{navy}}{}{0em}{}[\vspace{-2pt}\rule{\linewidth}{0.8pt}]
\titleformat{\subsection}[hang]{\normalsize\bfseries\color{navy}}{}{0em}{}
\titleformat{\subsubsection}[hang]{\normalsize\bfseries\color{darkgrey}}{}{0em}{}

% --- Header / Footer ---
\pagestyle{fancy}
\fancyhf{}
\fancyhead[L]{\small\color{midgrey}ठेकेदार हेमराज --- बचाव संक्षेप}
\fancyhead[R]{\small\color{midgrey}राजस्थान | निर्माण विधि}
\fancyfoot[C]{\small\color{midgrey}\thepage}
\fancyfoot[L]{\small\color{midgrey}गोपनीय --- केवल विधिक कार्यवाही हेतु}
\fancyfoot[R]{\small\color{midgrey}सभी IS उद्धरण BIS से सत्यापित करें}
\renewcommand{\headrulewidth}{0.4pt}
\renewcommand{\footrulewidth}{0.4pt}

% --- Frame styles ---
\newmdenv[linecolor=navy, linewidth=1.2pt, backgroundcolor=lightblue, skipabove=8pt, skipbelow=8pt, innertopmargin=8pt, innerbottommargin=8pt, innerleftmargin=10pt, innerrightmargin=10pt]{groundbox}

\newmdenv[linecolor=gold, linewidth=1.2pt, backgroundcolor=lightgold, skipabove=8pt, skipbelow=8pt, innertopmargin=8pt, innerbottommargin=8pt, innerleftmargin=10pt, innerrightmargin=10pt]{standardbox}

\newmdenv[linecolor=red!70, linewidth=1.2pt, backgroundcolor=lightred, skipabove=8pt, skipbelow=8pt, innertopmargin=8pt, innerbottommargin=8pt, innerleftmargin=10pt, innerrightmargin=10pt]{criticalbox}

\newmdenv[linecolor=navy!60, linewidth=0.8pt, backgroundcolor=lightgreen, skipabove=6pt, skipbelow=6pt, innertopmargin=6pt, innerbottommargin=6pt, innerleftmargin=8pt, innerrightmargin=8pt]{prayerbox}

% --- List settings ---
\setlist[itemize]{noitemsep, topsep=2pt, leftmargin=1.2em}
\setlist[enumerate]{noitemsep, topsep=2pt, leftmargin=1.5em}

% --- Spacing ---
\setstretch{1.4}

% ================================================================
\begin{document}
% ================================================================

% --- Title Page ---
\begin{titlepage}
  \centering
  \vspace*{1cm}
  {\Huge\bfseries\color{navy} ठेकेदार बचाव संक्षेप\par}
  \vspace{0.3cm}
  {\Large\color{gold} \rule{0.7\linewidth}{1pt}\par}
  \vspace{0.4cm}
  {\large\itshape राजय विधि विज्ञान प्रयोगशाला, जयपुर, राजस्थान की\\
  FSL रिपोर्ट को चुनौती\par}
  \vspace{0.8cm}

  \begin{mdframed}[linecolor=navy, linewidth=1.5pt, backgroundcolor=lightblue]
    \centering
    \textbf{विषय:}\\[6pt]
    \large\textbf{ठेकेदार हेमराज} (प्रतिवादी)\\
    \normalsize ग्राम/जिला, राजस्थान\\[10pt]
    \textbf{विरुद्ध:}\\[6pt]
    सीमा दीवार निर्माण स्थल से संग्रहीत निर्माण सामग्री\\
    (पैकेट A, B, C) पर FSL रिपोर्ट\\[10pt]
    \textbf{प्रयोगशाला:} राजय विधि विज्ञान प्रयोगशाला, जयपुर\\
    \textbf{हस्ताक्षरकर्ता:} KL वर्मा (SSO भौतिकी) एवं डॉ.\ शैलेन्द्र झा\\
    (सहायक निदेशक भौतिकी)
  \end{mdframed}

  \vspace{0.8cm}
  \textbf{परीक्षित नमूने:}

  \begin{center}
  \begin{tabular}{|l|l|l|}
    \hline
    \textbf{पैकेट} & \textbf{सामग्री} & \textbf{FSL परिणाम} \\
    \hline
    पैकेट A & सीमा दीवार पर सीमेंट प्लास्टर & सीमेंट : रेत = 1 : 18 \\
    पैकेट B & पत्थर चिनाई के जोड़ का मसाला & सीमेंट : रेत = 1 : 17 \\
    पैकेट C & नींव का कंक्रीट (पत्थर चिनाई दीवार) & सीमेंट : रेत : गिट्टी = 1 : 5.75 : 9.5 \\
    \hline
  \end{tabular}
  \end{center}

  \vspace{0.8cm}
  \begin{criticalbox}
    \centering
    \textbf{\large पाँच विशेषज्ञ एजेंटों द्वारा सत्यापित}\\[4pt]
    \small सभी भारतीय मानक संस्करण, धारा संख्या, अंतर्राष्ट्रीय मानक,\\
    वैज्ञानिक दावे एवं विधिक पूर्वनिर्णय स्वतंत्र रूप से सत्यापित\\
    2 सुधारे गये $\cdot$ 7 पुष्टि किये गये $\cdot$ 5 नये जोड़े गये $\cdot$ 1 सशक्त किया गया
  \end{criticalbox}

  \vfill
  {\small\color{midgrey}Legal Luminaire --- राजस्थान निर्माण विधि बचाव प्रणाली द्वारा तैयार\\
  \textbf{अस्वीकरण:} न्यायालय में उपयोग से पूर्व सभी IS उद्धरण BIS की मूल प्रकाशनों (www.bis.gov.in) से सत्यापित करें।\\
  तैयारी की तिथि: \today}
\end{titlepage}

% --- Table of Contents ---
\tableofcontents
\newpage

% ================================================================
\section*{प्रारंभिक टिप्पणी --- पाँच-एजेंट सत्यापन}
\addcontentsline{toc}{section}{प्रारंभिक टिप्पणी --- पाँच-एजेंट सत्यापन}
% ================================================================

इस संक्षेप की तैयारी से पूर्व, प्रत्येक उद्धरण, मानक संदर्भ, धारा संख्या, वैज्ञानिक दावे तथा विधिक पूर्वनिर्णय को पाँच विशेषज्ञ एजेंटों द्वारा स्वतंत्र रूप से सत्यापित किया गया:

\begin{itemize}
  \item \textbf{एजेंट 1:} समस्त IS मानक संस्करणों एवं BIS धारा संख्याओं की पुष्टि की। IS 1199:1959 को IS 1199 (Part 5):2018 द्वारा और IS 516:1959 को IS 516 (Part 1):2018 द्वारा प्रतिस्थापित पाया गया।
  \item \textbf{एजेंट 2:} ASTM C823:2017, ASTM C856:2020, BS EN 12504-1:2019 एवं BS EN 14630:2006 अंतर्राष्ट्रीय मानकों की धारा संख्या एवं वर्तमानता की पुष्टि की।
  \item \textbf{एजेंट 3:} वैज्ञानिक दावों की सटीकता की जाँच की: (क) कठोर प्लास्टर में कार्बोनेशन; (ख) पत्थर की सिलिका से मसाले का संदूषण; (ग) IS 4031 भाग 5 की कंक्रीट हेतु परिसीमा; (घ) IS 269:2015 के अनुसार सीमेंट सिलिका भिन्नता।
  \item \textbf{एजेंट 4:} सर्वोच्च न्यायालय के पूर्वनिर्णय उद्धरणों की पुष्टि की --- सभी सही पाये गये।
  \item \textbf{एजेंट 5:} हिन्दी अनुवाद की विधिक शुद्धता, शब्दावली एवं देवनागरी पाठ की समीक्षा की।
\end{itemize}

% ================================================================
\section{प्रस्तावना एवं सारांश}
% ================================================================

यह संक्षेप ठेकेदार हेमराज की ओर से, राजय विधि विज्ञान प्रयोगशाला, जयपुर द्वारा जारी FSL रिपोर्ट के उत्तर में प्रस्तुत किया जा रहा है। इस रिपोर्ट में यह दर्शाने का प्रयास किया गया है कि सीमा दीवार निर्माण में उपयोग की गई सामग्री घटिया थी। FSL द्वारा परीक्षित तीन नमूने थे:

\begin{itemize}
  \item \textbf{पैकेट A:} सीमा दीवार की बाहरी सतह पर सीमेंट प्लास्टर --- दावा अनुपात सीमेंट:रेत = 1:18
  \item \textbf{पैकेट B:} सीमा दीवार की पत्थर चिनाई के जोड़ों में सीमेंट मसाला --- दावा अनुपात सीमेंट:रेत = 1:17
  \item \textbf{पैकेट C:} पत्थर चिनाई दीवार की नींव का कंक्रीट --- दावा अनुपात सीमेंट:रेत:गिट्टी = 1:5.75:9.5
\end{itemize}

ये तीनों परिणाम (क) नौ प्रक्रियागत उल्लंघनों से और (ख) प्रत्येक नमूने के प्रकार के लिए विशिष्ट वैज्ञानिक एवं पद्धतिगत दोषों से \textbf{दूषित} हैं। इन दोषों के संयुक्त भार को देखते हुए, इस FSL रिपोर्ट के आधार पर ठेकेदार के विरुद्ध कोई प्रतिकूल निष्कर्ष नहीं निकाला जा सकता।

% ================================================================
\section{भाग-१ --- सीमा दीवार पर सीमेंट प्लास्टर (पैकेट A)}
% ================================================================

\subsection{पैकेट का विवरण}

\begin{center}
\begin{tabular}{|l|l|}
  \hline
  \textbf{सामग्री} & सीमा दीवार की बाहरी सतह पर सीमेंट प्लास्टर \\
  \textbf{FSL परिणाम} & सीमेंट : रेत = 1 : 18 \\
  \textbf{लागू विनिर्देश} & IS 1661:1972 --- बाहरी प्लास्टर हेतु CM 1:4 से CM 1:6 \\
  \textbf{मुख्य टिप्पणी} & 1:18 मानक श्रेणी से 300--450\% अधिक रेत दर्शाता है। \\
                         & खड़ी दीवार इस अनुपात को असंभव सिद्ध करती है। \\
  \hline
\end{tabular}
\end{center}

\subsection{नौ प्रक्रियागत आधार (पैकेट A पर लागू)}

\subsubsection{आधार 1 --- ठेकेदार के प्रतिनिधि की अनुपस्थिति में नमूना संग्रह \textbf{[अत्यंत महत्वपूर्ण]}}

\begin{groundbox}
प्लास्टर का नमूना ठेकेदार या उनके अधिकृत प्रतिनिधि को बिना किसी सूचना के और उनकी अनुपस्थिति में लिया गया। यह प्राकृतिक न्याय के सिद्धांत --- \textit{Audi Alteram Partem} --- का मूलभूत उल्लंघन है। ठेकेदार के गवाह के बिना इसकी कोई स्वतंत्र पुष्टि नहीं हो सकती: (क) किस दीवार की सतह से नमूना लिया गया; (ख) उपकरण स्वच्छ था या नहीं; (ग) एकत्रित सामग्री ठेकेदार के कार्य की थी या नहीं; (घ) नमूने को तुरंत लेबल किया और सील किया गया या नहीं।
\end{groundbox}

\begin{standardbox}
\textbf{IS 1199 (Part 5):2018, धारा 5:} नमूना संग्रह दोनों पक्षों के प्रतिनिधियों की उपस्थिति में हो। \textit{[IS 1199:1959 का स्थान लेता है।]}\\
\textbf{IS 516 (Part 1):2018, धारा 4:} विधिक कार्यवाही हेतु परिणामों में दोनों पक्षों को प्रतिनिधित्व का अधिकार है।\\
\textbf{विधिक सिद्धांत:} मेनका गांधी बनाम भारत संघ (1978) 1 SCC 248: किसी पक्ष को बाहर रखकर उत्पन्न साक्ष्य के आधार पर कोई प्रतिकूल आदेश नहीं दिया जा सकता।
\end{standardbox}

\subsubsection{आधार 2 --- तूफानी मौसम, भीगी सीढ़ी और असुरक्षित उपकरण \textbf{[अत्यंत महत्वपूर्ण]}}

\begin{groundbox}
तूफान और बारिश के दिन, जब जमीन पर लगभग दो फुट पानी भरा था, एक मजदूर भीगी स्टील की सीढ़ी पर छाता लगाकर खड़े होकर हथौड़े से दीवार पर प्रहार कर नमूना ले रहा था। 12--20 मिमी मोटे प्लास्टर पर वर्षाजल तुरंत प्रवेश करता है और उस सामग्री में उपस्थित कैल्शियम व सिलिका की मात्रा को बदल देता है जिसका बाद में रासायनिक विश्लेषण होता है। IS 1199 (Part 5):2018 और ASTM C823:2017 वर्षा में नमूना संग्रह स्पष्ट रूप से प्रतिबंधित करते हैं।
\end{groundbox}

\begin{standardbox}
\textbf{IS 1199 (Part 5):2018, धारा 5:} नमूनों को वर्षा एवं प्रतिकूल पर्यावरणीय परिस्थितियों से तुरंत संरक्षित किया जाए।\\
\textbf{BS EN 14630:2006, धारा 5.2:} वर्षा में नमूना नहीं लिया जाएगा; गीली जमीन पर गिरे नमूने त्यागे जाएँगे।\\
\textbf{ASTM C823/C823M:2017, धारा 8:} स्वच्छ तिरपाल, सुरक्षित एवं स्थिर पहुँच अनिवार्य।
\end{standardbox}

\subsubsection{आधार 3 --- अभिरक्षा श्रृंखला का टूटना \textbf{[अत्यंत महत्वपूर्ण]}}

\begin{groundbox}
प्लास्टर का नमूना उसी तूफान में बिना किसी प्रत्येक हस्तांतरण बिंदु पर हस्ताक्षरित अभिरक्षा-श्रृंखला दस्तावेज के परिवहन किया गया। IS/IEC 17025:2017 धारा 7.4 और NABL Doc 112:2018 धारा 5.8 यह दस्तावेजीकरण न्यूनतम आवश्यकता के रूप में अनिवार्य करते हैं। बिना अटूट अभिरक्षा श्रृंखला के यह स्थापित नहीं किया जा सकता कि जयपुर में जो परीक्षण हुआ वह ठेकेदार हेमराज की सीमा दीवार से आया था।
\end{groundbox}

\begin{standardbox}
\textbf{IS/IEC 17025:2017, धारा 7.4:} नमूनों के परिवहन, प्राप्ति, संभालने, भंडारण और निपटान की प्रलेखित प्रक्रिया आवश्यक।\\
\textbf{NABL Doc 112:2018, धारा 5.8:} प्रत्येक हस्तांतरण पर हस्ताक्षरित अभिरक्षा-श्रृंखला दस्तावेज।\\
\textbf{विधिक सिद्धांत:} मो.\ खालिद बनाम पश्चिम बंगाल राज्य (2002) 7 SCC 334: अभिरक्षा-श्रृंखला में अंतराल प्रतिवादी के पक्ष में उचित संदेह उत्पन्न करता है।
\end{standardbox}

\subsubsection{आधार 4 --- समय पर परीक्षण एवं भंडारण पर संदेह}

कठोर प्लास्टर कार्बोनेशन के प्रति अत्यंत संवेदनशील है --- Ca(OH)$_2$ + CO$_2$ $\rightarrow$ CaCO$_3$ + H$_2$O की प्रतिक्रिया से घुलनशील कैल्शियम घटता है और सीमेंट अंश कम दिखता है। वर्षा में एकत्रित और अनुचित ढंग से संग्रहीत प्लास्टर नमूने में परीक्षण से पूर्व तेज कार्बोनेशन होगी। FSL रिपोर्ट में न संग्रह की तारीख है, न प्राप्ति की, न परीक्षण की।

\subsubsection{आधार 5 --- संग्रह और परीक्षण के बीच समय सीमा का पालन नहीं}

FSL रिपोर्ट में न संग्रह तिथि, न प्राप्ति तिथि, न परीक्षण तिथि उल्लिखित है। IS 1199 (Part 5):2018 की धारा 6 के अनुसार यह रिपोर्ट का अनिवार्य भाग है। अनिवार्य परीक्षण समयसीमा का अनुपालन सत्यापित करने का भार जाँच प्राधिकारी पर है --- जो वे पूरा नहीं कर पाये हैं।

\subsubsection{आधार 6 --- ठेकेदार की अनुपस्थिति में परीक्षण \textbf{[अत्यंत महत्वपूर्ण]}}

\begin{groundbox}
1:18 का अनुपात देने वाला रासायनिक विश्लेषण ठेकेदार को बिना सूचना के और उनकी अनुपस्थिति में किया गया। NABL Doc 112:2018 धारा 5.10 के अनुसार विधिक विवाद में पक्षों को परीक्षण से पूर्व सूचना देना और गवाह नियुक्त करने का अधिकार देना अनिवार्य है। हस्ताक्षरकर्ता KL वर्मा का पदनाम SSO \textit{भौतिकी} है --- सीमेंटयुक्त सामग्री का रासायनिक विश्लेषण रसायन विज्ञान का कार्य है, भौतिकी का नहीं।
\end{groundbox}

\begin{standardbox}
\textbf{IS 516 (Part 1):2018, धारा 4:} विधिक कार्यवाही हेतु दोनों पक्षों को परीक्षण में उपस्थित रहने का अधिकार।\\
\textbf{NABL Doc 112:2018, धारा 5.10:} पक्षों को सूचना दी जाए; गवाह नियुक्त करने का अधिकार।\\
\textbf{विधिक सिद्धांत:} भारत संघ बनाम तुलसीराम पटेल (1985) 3 SCC 398: आरोपित व्यक्ति को प्रत्येक महत्वपूर्ण चरण में वास्तविक अवसर मिलना चाहिए।
\end{standardbox}

\subsubsection{आधार 7 --- स्वतंत्र पुन:परीक्षण हेतु नमूना संरक्षित नहीं}

पैकेट A का कोई भाग जज (रेफरी) नमूने के रूप में संरक्षित नहीं किया गया। IS 1199 (Part 5):2018 के अनुसार संग्रहीत मात्रा का कम से कम एक-तिहाई न्यूनतम 90 दिनों तक संरक्षित रखना अनिवार्य है। एक बार नमूना प्राथमिक परीक्षण में पूरी तरह उपयोग होने पर ठेकेदार का स्वतंत्र सत्यापन का अधिकार हमेशा के लिए समाप्त हो जाता है।

\subsubsection{आधार 8 --- FSL रिपोर्ट में कोई भारतीय मानक उद्धृत नहीं \textbf{[अत्यंत महत्वपूर्ण]}}

\begin{groundbox}
FSL रिपोर्ट में किसी भी भारतीय मानक का उल्लेख नहीं है जिसके अन्तर्गत रासायनिक विश्लेषण किया गया। IS/IEC 17025:2017 धारा 7.8.2 इसे प्रत्येक परीक्षण रिपोर्ट में अनिवार्य न्यूनतम आवश्यकता घोषित करती है। रिपोर्ट में स्वयं लिखा है: \textit{``नियंत्रण नमूना नहीं मिला, इसलिए यह मानकर कि उत्तम गुणवत्ता के सीमेंट में 21\% घुलनशील सिलिका होती है, अनुपात की गणना की गई।''} IS 4031 (Part 5):1988 धारा 4 के अनुसार सत्यापित सीमेंट नियंत्रण नमूना अनिवार्य है। इसके बिना परिणाम केवल संकेतात्मक है, निर्णायक नहीं।
\end{groundbox}

\begin{standardbox}
\textbf{IS 1661:1972 (Reaffirmed 2020), तालिका 1:} बाहरी प्लास्टर: CM 1:4 (सम्पन्न) से CM 1:6 (सामान्य)। 1:6 से अधिक कोई अनुपात किसी भी मान्यता प्राप्त भारतीय मानक में किसी भी कार्यात्मक बाहरी प्लास्टर के लिए नहीं है।\\
\textbf{IS 4031 (Part 5):1988 (Reaffirmed 2018), धारा 4:} रासायनिक विश्लेषण हेतु सीमेंट का नियंत्रण नमूना अनिवार्य।\\
\textbf{IS/IEC 17025:2017, धारा 7.8.2:} प्रत्येक परीक्षण रिपोर्ट में परीक्षण विधि संदर्भ अनिवार्य।
\end{standardbox}

\subsubsection{आधार 9 --- रिपोर्ट में परीक्षणों की संख्या का उल्लेख नहीं}

FSL रिपोर्ट एक एकल अनुपात प्रदान करती है बिना यह बताये कि यह एकल निर्धारण है या औसत। IS 4031 (Part 5):1988 धारा 13 और IS 1199 (Part 5):2018 न्यूनतम दो सुसंगत निर्धारणों की माँग करते हैं। कोई दोहरान डेटा, माप अनिश्चितता या व्यक्तिगत निर्धारण मूल्य नहीं दिया --- यह सब IS/IEC 17025:2017 धारा 7.8.3 के अन्तर्गत अनिवार्य हैं।

\subsection{सामग्री-विशिष्ट तकनीकी चुनौतियाँ (पैकेट A --- प्लास्टर)}

\subsubsection{चुनौती A1 --- कठोर प्लास्टर का कार्बोनेशन}

Ca(OH)$_2$ + CO$_2$ $\rightarrow$ CaCO$_3$ + H$_2$O की रासायनिक प्रतिक्रिया से सीमेंट का घुलनशील कैल्शियम अंश घटता है। IS 4031 (Part 5):1988 की रासायनिक गणना में यही अंश सीमेंट की मात्रा का आधार है। निर्माण के वर्षों बाद, वर्षाजल में भीगे और अनुचित ढंग से भंडारित नमूने में कार्बोनेशन से सीमेंट का अनुपात वास्तविकता से बहुत कम प्रतीत होगा। FSL का 1:18 का परिणाम आंशिक या पूर्ण रूप से कार्बोनेशन का परिणाम हो सकता है, मूल मिश्रण अनुपात का नहीं।

\subsubsection{चुनौती A2 --- आधारभूत सामग्री का संदूषण}

बाहरी प्लास्टर केवल 12--20 मिमी मोटा होता है। तूफान में भीगी सीढ़ी पर खड़े होकर हथौड़े से मारने पर अनिवार्यतः प्लास्टर की पतली परत के नीचे की पत्थर चिनाई के टुकड़े भी नमूने में आ जाते हैं। पत्थर के टुकड़ों में सिलिका की मात्रा बहुत अधिक होती है (क्वार्टज़ाइट: 92--95\% SiO$_2$)। इनकी सिलिका गणना में ``समुच्चय सिलिका'' के रूप में जुड़ जाती है, जो सीमेंट:रेत का अनुपात अत्यधिक रेत-युक्त दर्शाती है। FSL रिपोर्ट में किसी पेट्रोग्राफिक परीक्षण (ASTM C856:2020) का उल्लेख नहीं है जो यह सत्यापित करता।

\subsubsection{चुनौती A3 --- कार्यात्मक बाहरी प्लास्टर के लिए 1:18 की असंभाव्यता}

IS 1661:1972 एवं IS 2402:1963 बाहरी उपयोग हेतु 1:3 से 1:6 अनुपात निर्धारित करते हैं। 1:18 का प्लास्टर निर्माण के कुछ महीनों में ही झड़ जाता। खड़ी सीमा दीवार इस अनुपात की असंभाव्यता का प्रत्यक्ष प्रमाण है।

% ================================================================
\section{भाग-२ --- पत्थर चिनाई के जोड़ का मसाला (पैकेट B)}
% ================================================================

\subsection{पैकेट का विवरण}

\begin{center}
\begin{tabular}{|l|l|}
  \hline
  \textbf{सामग्री} & पत्थर चिनाई के जोड़ों में सीमेंट मसाला, सीमा दीवार \\
  \textbf{FSL परिणाम} & सीमेंट : रेत = 1 : 17 \\
  \textbf{लागू विनिर्देश} & IS 2250:1981 --- ग्रेड M3/M4 (1:5 अथवा 1:6) \\
  \textbf{मुख्य टिप्पणी} & 1:17 के आसपास कोई मान्यता प्राप्त चिनाई मसाला ग्रेड नहीं। \\
                         & 1:17 मसाले से बनी दीवार खड़ी नहीं रह सकती। \\
  \hline
\end{tabular}
\end{center}

\subsection{प्रक्रियागत आधार (पैकेट B पर लागू)}

पैकेट A के लिए कहे गए नौ प्रक्रियागत आधार (धारा 2.2) पैकेट B पर समान रूप से लागू हैं। विशेष अनुकूलन:

\textbf{आधार 2 (अनुकूलित):} पत्थर चिनाई के जोड़ मात्र 10--25 मिमी चौड़े होते हैं। वर्षा में भीगी पत्थर चिनाई से हथौड़े द्वारा बिना पत्थर के टुकड़े शामिल किये शुद्ध मसाला निकालना शारीरिक रूप से असंभव है। जोड़ों में भरा वर्षाजल नमूने के गिरने से पहले ही उसे दूषित कर देता है।

शेष आधार (3--9) पैकेट A के समान ही लागू होते हैं।

\subsection{सामग्री-विशिष्ट तकनीकी चुनौतियाँ (पैकेट B --- चिनाई मसाला)}

\subsubsection{चुनौती B1 --- पत्थर की सिलिका से संदूषण: निर्णायक वैज्ञानिक दोष \textbf{[सर्वाधिक महत्वपूर्ण]}}

\begin{criticalbox}
\textbf{यह संपूर्ण संक्षेप में सर्वाधिक महत्वपूर्ण वैज्ञानिक तर्क है।}

IS 4031 भाग 5:1988 की रासायनिक विश्लेषण विधि मानती है कि नमूने में सम्पूर्ण सिलिका केवल सीमेंट और रेत से आती है। पत्थर चिनाई के जोड़ों से शुद्ध मसाला हथौड़े से नहीं निकाला जा सकता --- पत्थर के टुकड़े अनिवार्यतः नमूने में आते हैं। राजस्थान के पत्थर:
\begin{itemize}
  \item क्वार्टज़ाइट: 92--95\% SiO$_2$
  \item बलुआ पत्थर: 60--80\% SiO$_2$
  \item ग्रेनाइट: 60--75\% SiO$_2$
\end{itemize}

इन पत्थर के टुकड़ों की सिलिका गणना में ``समुच्चय सिलिका'' बन जाती है। नमूने में मात्र 5\% पत्थर टुकड़ा (जो आँखों को दिखे भी न) 1:6 के वास्तविक मसाले को 1:14 या इससे भी दुर्बल दर्शा सकता है। पत्थर संदूषण + कार्बोनेशन + 21\% सिलिका की अप्रमाणित मान्यता --- ये तीनों एक ही दिशा में कार्य करते हुए मसाले को वास्तविकता से अत्यधिक दुर्बल दर्शाते हैं --- और इस प्रकार 1:17 के परिणाम की पूरी तरह व्याख्या बिना किसी वास्तविक दोष के हो जाती है।
\end{criticalbox}

\begin{standardbox}
\textbf{IS 2250:1981 (Reaffirmed 2020), धारा 6.1:} मसाले के नमूने केवल जोड़ की छेनी से निकाले जाएँ, चिनाई इकाई की सामग्री शामिल न हो।\\
\textbf{IS 1597 (Part 1):1992 (Reaffirmed 2018), धारा 6:} पत्थर चिनाई से मसाला परीक्षण में मसाले को पत्थर से अलग करना अनिवार्य।\\
\textbf{ASTM C856:2020, धारा 9:} पेट्रोग्राफिक परीक्षण ही मसाले को पत्थर टुकड़ों से पृथक पहचानने की मानक विधि है।
\end{standardbox}

\subsubsection{चुनौती B2 --- 1:17 की किसी कार्यात्मक चिनाई मसाले के लिए असंभाव्यता}

IS 2250:1981 तालिका 1 में M1 (1:3) से M5 (1:8) तक ग्रेड हैं --- 1:17 के आसपास कोई ग्रेड नहीं। 1:17 के मसाले में कोई संपीड़न सामर्थ्य नहीं होगी और दीवार अपने भार से ही ढह जाती। खड़ी पत्थर चिनाई सीमा दीवार स्वयं सिद्ध करती है कि प्रयुक्त मसाला पर्याप्त मजबूत था।

\subsubsection{चुनौती B3 --- 21\% सिलिका मान्यता और उसका संचयी प्रभाव}

IS 269:2015 (OPC विनिर्देश) सीमेंट में केवल न्यूनतम 17\% SiO$_2$ की आवश्यकता है --- कोई ऊपरी सीमा नहीं; वास्तविक भिन्नता लगभग 19--24\% है। FSL की 21\% मान्यता किसी भारतीय मानक द्वारा अनुमोदित नहीं है। पैकेट B में यह त्रुटि पत्थर संदूषण त्रुटि के साथ संचित होती है --- दोनों मिलकर मिश्रण को दुर्बल दर्शाती हैं।

% ================================================================
\section{भाग-३ --- नींव का कंक्रीट (पैकेट C)}
% ================================================================

\subsection{पैकेट का विवरण}

\begin{center}
\begin{tabular}{|l|l|}
  \hline
  \textbf{सामग्री} & पत्थर चिनाई दीवार की नींव का कंक्रीट \\
  \textbf{FSL परिणाम} & सीमेंट : रेत : गिट्टी = 1 : 5.75 : 9.5 \\
  \textbf{लागू मानक} & IS 456:2000 (Reaffirmed 2021) --- नींव हेतु M10 या M15 \\
  \textbf{निर्णायक दोष} & IS 4031 (भाग 5) \textbf{सीमेंट परीक्षण मानक} है, \\
                        & \textbf{कंक्रीट विश्लेषण मानक नहीं।}\\
                        & FSL ने \textbf{गलत विधि} लागू की। \\
  \hline
\end{tabular}
\end{center}

\subsection{प्रक्रियागत आधार (पैकेट C पर लागू)}

पैकेट A के नौ आधार पैकेट C पर समान रूप से लागू हैं।

\textbf{आधार 8 (अतिरिक्त अनुकूलन):} पैकेट C के लिए न केवल रिपोर्ट में IS का उल्लेख नहीं है, बल्कि जो विधि उपयोग की गई (IS 4031 भाग 5 रासायनिक विश्लेषण) वह IS 4031 (Part 5):1988 के दायरे से बाहर है। मौजूदा संरचनाओं में कंक्रीट सत्यापन की सही विधियाँ हैं: IS 516 (Part 5):2020 (कोर ड्रिलिंग) और ASTM C856:2020 (पेट्रोग्राफिक परीक्षण)।

\subsection{सामग्री-विशिष्ट तकनीकी चुनौतियाँ (पैकेट C --- नींव कंक्रीट)}

\subsubsection{चुनौती C1 --- मौलिक रूप से गलत विश्लेषण विधि \textbf{[निर्णायक]}}

\begin{criticalbox}
\textbf{यह संपूर्ण संक्षेप का एकल सर्वाधिक महत्वपूर्ण तकनीकी तर्क है।}

IS 4031 भाग 5:1988 धारा 1 के अनुसार यह विधि \textbf{सीमेंट के विश्लेषण} तथा मसाले (सीमेंट + बारीक समुच्चय) के लिए है। यह \textbf{कंक्रीट के लिए नहीं है} जिसमें मोटा समुच्चय (गिट्टी/पत्थर चिप्स) भी होता है।

कंक्रीट में मोटे समुच्चय में स्वयं सिलिका होती है:
\begin{itemize}
  \item क्वार्टज़ाइट गिट्टी: 92--95\% SiO$_2$
  \item कुचला ग्रेनाइट: 60--75\% SiO$_2$
\end{itemize}

IS 4031 भाग 5 की गणना में ``अघुलनशील अवशेष'' से कुल समुच्चय निकाला जाता है। मसाले में यह केवल रेत है; कंक्रीट में इसमें रेत और गिट्टी दोनों होती हैं और इन्हें अलग नहीं किया जा सकता। तीन-घटक अनुपात (1:5.75:9.5) इस विधि से निकालने का प्रयास वैज्ञानिक रूप से अस्वीकार्य है।

सही विधियाँ जो FSL ने नहीं अपनाईं:
\begin{itemize}
  \item \textbf{IS 516 (Part 5):2020} --- कोर ड्रिलिंग और संपीड़न परीक्षण
  \item \textbf{ASTM C856:2020} --- कठोर कंक्रीट का पेट्रोग्राफिक परीक्षण
  \item \textbf{IS 13311 (Part 1)\&(2):1992} --- अविनाशी परीक्षण (UPV एवं रिबाउंड हैमर)
  \item \textbf{BS EN 12504-1:2019} --- मौजूदा संरचनाओं से कोर निकालना
\end{itemize}
\end{criticalbox}

\begin{standardbox}
\textbf{IS 456:2000 (Reaffirmed 2021), धारा 14--15:} कंक्रीट की गुणवत्ता कोर की संपीड़न सामर्थ्य परीक्षण से सत्यापित होती है।\\
\textbf{IS 516 (Part 5):2020:} मौजूदा कंक्रीट संरचनाओं में कोर ड्रिलिंग अनिवार्य विधि है।\\
\textbf{IS 4031 (Part 5):1988 (Reaffirmed 2018), धारा 1 --- दायरा:} यह विधि हाइड्रॉलिक सीमेंट के लिए है; कंक्रीट पर इसका प्रयोग दायरे से बाहर है।\\
\textbf{ASTM C856:2020:} पेट्रोग्राफिक परीक्षण रेत और मोटे समुच्चय को अलग पहचान सकता है --- रासायनिक विश्लेषण नहीं।
\end{standardbox}

\subsubsection{चुनौती C2 --- मोटे समुच्चय की सिलिका गणना को दूषित करती है}

FSL रिपोर्ट में S1 (रेत), S2 (गिट्टी), S3 (गिट्टी नमूना) प्राप्त होना दर्ज है लेकिन किसी का भी विश्लेषण परिणाम नहीं दिया गया। तीन-घटक अनुपात की गणना के लिए प्रत्येक समुच्चय की सिलिका सामग्री स्वतंत्र रूप से ज्ञात होना आवश्यक है। बिना S2/S3 के परिणाम के 1:5.75:9.5 का अनुपात एक अपारदर्शी, असत्यापनीय गणना है।

\subsubsection{चुनौती C3 --- विनिर्देश संदर्भ और भौतिक साक्ष्य}

IS 456:2000 तालिका 5 में सौम्य परिस्थितियों में सादे कंक्रीट के लिए नींव में न्यूनतम M10 ग्रेड निर्धारित है। कंक्रीट का विनिर्देश \textbf{सामर्थ्य ग्रेड} में होता है, आयतन अनुपात में नहीं। आयतन अनुपात की तुलना सामर्थ्य-आधारित विनिर्देश से करने के लिए अनेक अप्रमाणित मान्यताएँ आवश्यक हैं। इसके अतिरिक्त, नींव कंक्रीट निर्माण के बाद से पत्थर चिनाई दीवार का भार वहन कर रहा है --- यह इसकी पर्याप्त गुणवत्ता का प्रत्यक्ष प्रमाण है।

% ================================================================
\section{21\% घुलनशील सिलिका मान्यता --- तीनों पैकेट को प्रभावित करती है}
% ================================================================

\begin{criticalbox}
FSL रिपोर्ट में स्पष्ट लिखा है: \textit{``नियंत्रण नमूना नहीं मिला, इसलिए यह मानकर कि उत्तम गुणवत्ता के सीमेंट में 21\% घुलनशील सिलिका होती है, अनुपात की गणना की गई।''}

यह एकल वाक्य तीनों परिणामों की मात्रात्मक आधार को एक साथ कमज़ोर करता है:
\begin{enumerate}
  \item IS 269:2015 (OPC) केवल न्यूनतम 17\% SiO$_2$ की आवश्यकता है --- कोई ऊपरी सीमा नहीं। वास्तविक भिन्नता लगभग 19--24\%।
  \item 2\% की मान्यता त्रुटि अनुपात गणना में लगभग 10\% की त्रुटि उत्पन्न करती है।
  \item पैकेट A (कार्बोनेशन), पैकेट B (पत्थर संदूषण), पैकेट C (गलत विधि) के विशिष्ट दोष इस त्रुटि के साथ संचित होकर अत्यंत असंगत परिणाम देते हैं।
  \item IS 4031 (Part 5):1988 धारा 4 सत्यापित सीमेंट नियंत्रण नमूने को अनिवार्य करती है। जहाँ अनुपलब्ध हो, परिणाम को ``केवल संकेतात्मक'' घोषित करना अनिवार्य है --- FSL रिपोर्ट में यह अनिवार्य योग्यता नहीं है।
\end{enumerate}
\end{criticalbox}

% ================================================================
\section{निष्कर्ष एवं प्रार्थना}
% ================================================================

FSL रिपोर्ट अनेक स्वतंत्र आधारों पर दूषित है, जिनमें से कोई एक भी साक्ष्य को बाहर करने या न्यूनतम महत्व देने के लिए पर्याप्त होगा:

\begin{itemize}
  \item \textbf{नौ प्रक्रियागत उल्लंघन} तीनों पैकेट पर लागू: ठेकेदार प्रतिनिधि की अनुपस्थिति; तूफानी परिस्थितियों में संग्रह; अभिरक्षा-श्रृंखला टूटना; भंडारण एवं परीक्षण समयरेखा असत्यापित; नमूना संरक्षित नहीं; कोई IS उद्धृत नहीं; परीक्षणों की संख्या नहीं बताई।
  \item \textbf{सामग्री-विशिष्ट वैज्ञानिक दोष:} पैकेट A के लिए कार्बोनेशन और आधारभूत संदूषण; पैकेट B के लिए पत्थर सिलिका संदूषण; पैकेट C के लिए मौलिक रूप से गलत विधि।
  \item \textbf{21\% सिलिका की अप्रमाणित मान्यता} तीनों परिणामों की मात्रात्मक आधार को कमज़ोर करती है।
  \item \textbf{कोई सीमेंट नियंत्रण नमूना नहीं} --- IS 4031 (Part 5):1988 धारा 4 के अन्तर्गत सभी परिणाम केवल संकेतात्मक।
\end{itemize}

\begin{prayerbox}
\textbf{प्रार्थना} --- यह विनम्रतापूर्वक प्रार्थना की जाती है कि:
\begin{enumerate}
  \item पैकेट A, B एवं C के FSL परिणामों को विचार से बाहर किया जाए, अथवा कोई प्रमाणात्मक महत्व न दिया जाए, इस आधार पर कि वे नमूना संग्रह, अभिरक्षा-श्रृंखला, परीक्षण और रिपोर्टिंग के लागू भारतीय मानकों के उल्लंघन में तैयार किये गये।
  \item वैकल्पिक रूप से, एक नया परीक्षण निर्देशित किया जाए: (क) ठेकेदार को कम से कम 48 घंटे पूर्व सूचना; (ख) ठेकेदार के प्रतिनिधि की उपस्थिति में संग्रह एवं परीक्षण; (ग) प्रत्येक सामग्री की सही विधि से: प्लास्टर के लिए IS 1661/IS 516, मसाले के लिए IS 2250, कंक्रीट के लिए IS 516 (Part 5):2020 कोर ड्रिलिंग; (घ) सत्यापित सीमेंट नियंत्रण नमूने एवं मूल निर्माण में प्रयुक्त वास्तविक समुच्चय नमूनों सहित।
  \item ठेकेदार को सीमेंट बैच रिकॉर्ड, समुच्चय गुणवत्ता प्रमाण-पत्र और मिश्रण डिजाइन दस्तावेज़ीकरण आगे साक्ष्य में प्रस्तुत करने का अधिकार दिया जाए।
  \item मौजूदा FSL रिपोर्ट के आधार पर कोई प्रतिकूल आदेश नहीं दिया जाए।
\end{enumerate}
\end{prayerbox}

\vspace{1cm}
\begin{mdframed}[linecolor=red!60, linewidth=1pt, backgroundcolor=lightred]
\textbf{महत्वपूर्ण अस्वीकरण:} इस दस्तावेज़ में दिये गये सभी भारतीय मानक उद्धरणों को किसी भी विधिक या अर्ध-विधिक कार्यवाही में उपयोग से पूर्व भारतीय मानक ब्यूरो (BIS) की मूल प्रकाशनों --- \texttt{www.bis.gov.in} --- से स्वतंत्र रूप से सत्यापित किया जाना चाहिए। यह दस्तावेज़ सर्वोत्तम उपलब्ध जानकारी से तैयार एक विधिक सहायक साधन है, न कि सत्यापित विधिक परामर्श।
\end{mdframed}

\end{document}
